% Front matter of the Setting Up volume: edition notice (PUB-REF-0001 AN-20) and About This Book
% (AN-21 F-LOT: Who Should Read This Book / Conventions / Organization /
% Other Initech 123 Documentation / After You Complete This Book; annex
% Sec 5).  All boilerplate newly written (VO-12).
% GATES: test-pubs-build
\Gates{test-pubs-build}
\begin{TitleVerso}
\IsThisGoodForTheCompany\par
{\imLegalFont
\textbf{The following paragraph does not apply to any country where such
provisions are inconsistent with local law:} INITECH SYSTEMS CORPORATION
PROVIDES THIS PUBLICATION ``AS IS'' WITHOUT WARRANTY OF ANY KIND, EITHER
EXPRESS OR IMPLIED, INCLUDING, BUT NOT LIMITED TO, THE IMPLIED WARRANTIES OF
MERCHANTABILITY AND FITNESS FOR A PARTICULAR PURPOSE.  Some states do not
allow the exclusion of implied warranties, so the above exclusion may not
apply to you.\par
This publication may contain technical inaccuracies or typographical errors.
Its contents are revised from time to time, and the revisions appear in new
editions.  Initech Systems Corporation may change the products and programs
described here at any time without notice.\par
Copies of Initech publications, and technical information about Initech
products, are available from your authorized Initech dealer or your Initech
representative.\par}

{\footnotesize
Comments on this book are welcome; address them to: \InitechAddressLine.
\draft{PUB-REF-0001 O-5 and D-12: the address is to be checked against real
businesses before print}
Initech may use or distribute the information you supply in any way it
believes appropriate without incurring any obligation to you.\par
The software described in this publication is furnished under the Initech
Program License Agreement and may be used or copied only as that agreement
permits.\par}
\end{TitleVerso}

\FrontChapter{About This Book}
This volume contains four books.  \doctitle{Setting Up Initech 123} tells you
what comes in the package, what your computer needs, and how to start and
end Initech 123.  The \doctitle{Tutorial} teaches you Initech 123 in a series
of lessons, each one building on the one before.  \doctitle{Quick Start} is
a fast tour for people who already know worksheet programs.  \doctitle{Sample
Applications} gives you ready-to-type worksheets for common business tasks.

\section{Who Should Read This Book}
Everyone should read \doctitle{Setting Up Initech 123} before using Initech
123 for the first time.  If you have not used a worksheet program before,
work through the \doctitle{Tutorial} next, lesson by lesson, at the computer.
If you have, \doctitle{Quick Start} shows you what you need in an hour.  Use
the \doctitle{Sample Applications} when you are ready for a complete
business worksheet.

\section{Conventions}
This volume uses the following conventions.
\begin{itemize}
\item Names of keys are in small capitals.  A function key that has a name
  in Initech 123 is given by its name followed by the key in parentheses.\\
  Example: Press \fkey{Calc}{F9} to recalculate the worksheet.
\item Two key names joined by a hyphen mean: hold down the first key and
  press the second.\\
  Example: Press \kname{CTRL-BREAK} to return to \ind{READY} mode.
\item Text that you type is in this typeface.\\
  Example: Type \type{Paper} and press \kname{ENTER}.
\item In the lessons, each step is a line that begins with what you do, in
  bold, and continues with what to type or press, often followed by its
  purpose in italics.\\
  Example: \textbf{Press}\quad\kname{HOME} \textit{to move to cell A1}
\item Commands are named by the menu words you select, after a slash, as
  they appear in the control panel.\\
  Example: \slashcmd{/File Save}
\item Words in \textbf{bold} in the text are defined where they first
  appear; the \doctitle{Initech 123 Reference} has a glossary of them.
\end{itemize}

\section{Organization of This Book}
\begin{LotTable}{p{40mm}p{88mm}}
\LotHead{\textbf{Part} & \textbf{Contents}}
\doctitle{Setting Up Initech 123} (SU) & The package, the computer, starting
  and ending Initech 123, the keyboard, your worksheet files, problems\\
\doctitle{Tutorial} (TU) & Twenty-eight lessons in seven chapters\\
\doctitle{Quick Start} (QS) & Five short chapters, each ending with a Task
  Summary\\
\doctitle{Sample Applications} (SA) & Seven business worksheets\\
Master Index & One index for the four books; each locator line begins with
  the code of its book\LotEnd
\end{LotTable}

\section{Other Initech 123 Documentation}
\begin{LotTable}{p{50mm}p{78mm}}
\LotHead{\textbf{Manual} & \textbf{Description}}
\doctitle{Setting Up Initech 123, Tutorial, Quick Start and Sample
  Applications} (1051-0000-11) & This volume\\
\doctitle{Initech 123 Reference} (1051-0000-10) & Every command, @function
  and key, in full\\
\doctitle{Initech 123 Quick Reference} (1051-0000-20) & The keys, the
  @functions and the menus in a pocket booklet\\
Keyboard templates (1051-0000-21, 1051-0000-22) & The function-key
  assignments, for the 101-key and the 84-key keyboard\LotEnd
\end{LotTable}

\section{After You Complete This Book}
When you have finished the \doctitle{Tutorial} or \doctitle{Quick Start},
keep the \doctitle{Initech 123 Reference} at hand: it describes every command
in the order of the menus.
